\begin{center}
\textsc{Gauss} an \textsc{Gerling}.\quad Göttingen, 21.~Februar 1822.
\end{center}

\par\noindent\hspace*{\parindent}\mbox{.\ .\ .\ .\ .\ .}\quad Wenn der Anschluss an die Göttinger Sternwarte und der an den Inselsberg nicht auf Einem Platze auf dem Meisner geschehen kann, so dürfen Sie nun \textls[120]{beide} Plätze anwenden.
Es ist dabei gar nicht viel Vermehrung der Arbeit; denn nach den Principien der Wahrscheinlichkeitsrechnung brauchen Sie die Winkel zwischen den Objecten, die an beiden Punkten sichtbar sind, insofern Sie sie an beiden messen, an jedem nur halb so oft zu repetiren, als Sie gethan haben würden, wenn Sie nur Einen Standpunkt gebraucht hätten.
Die gegenseitige Lage beider Punkte gegen einander mit der grössten Schärfe und doch ohne überflüssige Mühe zu erhalten, wird Ihnen nicht schwer \mbox{fallen.\quad.\ .\ .\ .\ .\ .}
